\documentclass[11pt]{article}
\usepackage{fontspec,amsmath,amssymb,amsthm,hyperref}
\usepackage[margin=1in]{geometry}
\setlength{\emergencystretch}{3em}
\title{KIPLMC1 reuses the shared exponential integrator}
\author{LangevinLean mathematical exposition}
\date{}
\begin{document}\maketitle
\noindent Source module: \texttt{KIPLMC/Exponential.lean}. Compile this proof document with XeLaTeX or LuaLaTeX. The accompanying HTML pairs the mathematical argument with the exact Lean statements and records the scope of the formalization.


There is one exponential-step implementation: \texttt{KineticLangevin.exponentialStep}. KIPLMC1 supplies the proved averaged particle force and the common scaled-coordinate thermal amplitude \(\sqrt{2\gamma/N}\). The coordinate identities recover the paper's Algorithm 1 directly, including its separate velocity and position innovations. Noise covariance and independent-Gaussian law conditions must still be established for any stochastic convergence theorem.


\section{Supply the particle temperature to the shared map}


\texttt{kineticAmplitude} is \(a_N=\sqrt{2\gamma/N}\). \texttt{kiplmc1Step} calls the shared exponential map with force \(F\) and additive phase innovation \(a_N\,\text{noise}\). Consequently every block uses exactly the same friction response functions. In physical time and scaled coordinates this is an inverse temperature \(N\) kinetic system. The definition accepts arbitrary innovations; when used as the algorithm their velocity–position pair must have the common Gram covariance from the kinetic library.



\noindent Lean declarations: \texttt{kineticAmplitude}, \texttt{kiplmc1Step}.

\section{Recover both parameter updates}


Parameter coordinates are unchanged by spatial rescaling. Substitute \(F_\theta=N^{-1}\sum_i\nabla_\theta U_i\) in the two shared expressions:

\[V^{\theta+}=\psi_0V^\theta-\frac{\psi_1}N\sum_i\nabla_\theta U_i+a_N\xi_\theta,\quad
\theta^+=\theta+\psi_1V^\theta-\frac{\psi_2}N\sum_i\nabla_\theta U_i+a_N\xi'_\theta.\]

The proofs unfold the shared map, evaluate scalar multiplication in Euclidean coordinates, rewrite the proved force identity, and normalize division as multiplication by an inverse. The distinct projections of \texttt{noise} retain the correlated but generally unequal velocity and position increments.



\noindent Lean declarations: \texttt{kiplmc1Step\_parameter\_velocity}, \texttt{kiplmc1Step\_parameter\_position}.

\section{Undo the latent temperature and force scaling}


For \(N>0\), \(\sqrt N\sqrt{2\gamma/N}=\sqrt{2\gamma}\). Lean combines square roots using nonnegativity of \(N\), then cancels the positive particle number inside the root. Multiply each latent coordinate of the shared step by \(\sqrt N\). The force identity gives \(\sqrt NF_i=\nabla_xU_i\); the amplitude identity restores \(\sqrt{2\gamma}\) noise. Thus the physical latent velocity and position updates have the unaveraged pair force and the paper's latent thermal scale. The two coordinate proofs distribute multiplication, rewrite these exact identities and finish scalar rearrangement with \texttt{ring}.



\noindent Lean declarations: \texttt{unscale\_kineticAmplitude}, \texttt{kiplmc1Step\_particle\_velocity}, \texttt{kiplmc1Step\_particle\_position}.

\section{Construct the numerical recursion from a supplied innovation stream}


\texttt{kiplmc1Chain} starts from the given phase and uses innovation number \(n\) in the successor update. Each step is the shared exponential map specialized to the actual particle force. Recursive construction proves there is an explicit sequence for every prescribed stream; it is not an assumed sequence satisfying a desired estimate. To obtain the paper's actual random chain, one must additionally equip this stream with independent centered Gaussian pairs of the right covariance and prove its independence from initialization, measurability and finite moments.



\noindent Lean declarations: \texttt{kiplmc1Chain}.

\end{document}
